
\chapter{伪装导数题}
\section{概率}
\begin{example}{}{}
事件 $A$、$B$ 是随机试验中的两个事件，$P(A)$、$P(B) \neq 0$，且 $P(A)+P(B)=\frac{1}{2}$，\[\frac{1}{P(B|A)}+\frac{1}{P(A|B)}=\frac{1}{P(B)-P(AB)}+\frac{1}{P(A)-P(AB)}+2\]下列说法正确的是\\
\begin{tabular}{@{}p{0.45\textwidth} p{0.45\textwidth}@{}}
A. 事件 $A$、$B$ 一定不相互独立 & B. 若 $P(A)=\frac{1}{4}$，则 $P(AB)=\frac{1}{16}$ \\
C. $P(AB)\in\left(0,\frac{3-2\sqrt{2}}{4}\right]$ & D. $P(AB)\in\left(0,\frac{\sqrt{17}-1}{4}\right]$
\end{tabular}
\end{example}
\begin{solution}
设 $P(A) = x$，$P(B) = y$，$P(AB) = z$。则 $x > 0$，$y > 0$ 且 $x+y = \frac{1}{2}$。条件概率 $P(B|A) = \frac{P(AB)}{P(A)} = \frac{z}{x}$， $P(A|B) = \frac{P(AB)}{P(B)} = \frac{z}{y}$。
由于原等式中包含 $\frac{1}{P(B|A)}$ 等项作分母，可知 $z \neq 0$，即 $z > 0$。同时分母 $P(A)-P(AB) = x-z \neq 0$ 且 $P(B)-P(AB) = y-z \neq 0$，说明 $z < x$ 且 $z < y$。结合 $x+y=\frac{1}{2}$，必然有 $z < \frac{1}{4}$。将上述变量代入原等式：
\[\frac{x}{z} + \frac{y}{z} = \frac{1}{2z}=\frac{1}{y-z} + \frac{1}{x-z} + 2\Leftrightarrow xy = \frac{3}{2}z - z^2\]
已知 $x+y = \frac{1}{2}$ 且 $xy = \frac{3}{2}z - z^2$。根据韦达定理，$x$ 和 $y$ 是关于 $t$ 的一元二次方程 $t^2 - \frac{1}{2}t + \left(\frac{3}{2}z - z^2\right) = 0$ 的两个实数根。
为了保证实数根存在，判别式必须满足 $\Delta \ge 0$，即 $16z^2 - 24z + 1 \ge 0$。
解得 $z \le \frac{3-2\sqrt{2}}{4}$ 或 $z \ge \frac{3+2\sqrt{2}}{4}$。
结合前提 $0 < z < \frac{1}{4}$，且 $\frac{3-2\sqrt{2}}{4} \approx 0.0428 < \frac{1}{4}$，可得 $z$ 的严格取值范围为：
\[z \in \left(0, \frac{3-2\sqrt{2}}{4}\right]\]
\begin{itemize}
    \item \textbf{选项 A：}假设事件 $A, B$ 相互独立，则 $P(AB) = P(A)P(B)$，即 $z = xy$。代入推导出的关系式 $xy = \frac{3}{2}z - z^2$，得 $z = \frac{3}{2}z - z^2 \implies z^2 - \frac{1}{2}z = 0$。解得 $z=0$ 或 $z=\frac{1}{2}$。但这都不在有效定义域 $\left(0, \frac{3-2\sqrt{2}}{4}\right]$ 内。因此假设不成立，事件 $A, B$ 一定不相互独立，\textbf{故 A 正确}。
    \item \textbf{选项 B：}若 $P(A) = \frac{1}{4}$，由 $P(A)+P(B)=\frac{1}{2}$ 得 $P(B) = \frac{1}{4}$。此时 $xy = \frac{1}{16}$。代入关系式有 $\frac{1}{16} = \frac{3}{2}z - z^2$，化简为 $16z^2 - 24z + 1 = 0$。解得 $z = \frac{3-2\sqrt{2}}{4}$（另一根大于 $\frac{1}{4}$ 故舍去）。因为 $\frac{3-2\sqrt{2}}{4} \neq \frac{1}{16}$，\textbf{故 B 错误}。
    \item \textbf{选项 C：}$P(AB) = z$ 的取值范围确为 $\left(0, \frac{3-2\sqrt{2}}{4}\right]$，\textbf{故 C 正确}。
    \item \textbf{选项 D：}由于C正确，\textbf{故 D 正确}。
\end{itemize}
综上所述，正确的说法是 ACD。
\end{solution}


\begin{example}{}{}
已知无穷数列$\left\{a_n\right\}$满足$a_n\in\mathbf{N}^*$,且$\left\{a_n\right\}$前$n$项的中位数为$n.$\\
(1)若$\{a_n\}$是递增数列，求$\{a_n\}$的通项公式；\\
(2)证明：$a_n\geqslant n$ ;\\
(3)给定$k\in\mathbb{N}^*$,求关于$n$的方程$a_n=k$所有可能的解集.
\end{example}
\begin{solution}
(1) 由于 $\{a_n\}$ 是递增数列，当 $n=2k-1$ 时，前 $2k-1$ 项的中位数为第 $k$ 项，故 $a_k = 2k-1$；当 $n=2k$ 时，前 $2k$ 项的中位数为 $\frac{a_k + a_{k+1}}{2} = 2k$。将 $a_k = 2k-1$ 代入解得 $a_{k+1} = 2k+1$。故 $\{a_n\}$ 的通项公式为 $a_n = 2n-1$。

(2) 记 $S_n = \{a_1, a_2, \dots, a_n\}$，将其元素升序排列为 $x_{n, 1} \leqslant x_{n, 2} \leqslant \dots \leqslant x_{n, n}$，并记 $C_n(x)$ 为 $S_n$ 中不超过 $x$ 的元素个数。
对于 $n=2k-1$，中位数为 $x_{2k-1, k} = 2k-1$，说明 $C_{2k-1}(2k-2) \leqslant k-1$ 且 $C_{2k-1}(2k-1) \geqslant k$。
对于 $n=2k$，中位数满足 $x_{2k, k} + x_{2k, k+1} = 4k$。因为 $x_{2k, k} \leqslant 2k \leqslant x_{2k, k+1}$，无论 $x_{2k, k} = 2k$ 还是 $x_{2k, k} \leqslant 2k-1$，均有 $C_{2k}(2k-1) \leqslant k$。
由于 $S_{2k} = S_{2k-1} \cup \{a_{2k}\}$，故 $k \geqslant C_{2k}(2k-1) = C_{2k-1}(2k-1) + I(a_{2k} \leqslant 2k-1) \geqslant k + I(a_{2k} \leqslant 2k-1)$，这迫使 $I(a_{2k} \leqslant 2k-1) = 0$，即 $a_{2k} \geqslant 2k$。
同理，对于 $n=2k-2\ (k \geqslant 2)$，中位数满足 $x_{2k-2, k-1} + x_{2k-2, k} = 4k-4$。若 $x_{2k-2, k} \leqslant 2k-2$，则必有 $x_{2k-2, k-1} = x_{2k-2, k} = 2k-2$，导致 $C_{2k-2}(2k-2) \geqslant k$。但 $S_{2k-1}$ 包含 $S_{2k-2}$，这会导致 $k-1 \geqslant C_{2k-1}(2k-2) \geqslant C_{2k-2}(2k-2) \geqslant k$ 的矛盾。
故 $x_{2k-2, k} \geqslant 2k-1$，使得 $C_{2k-2}(2k-2) = k-1$。进而 $k-1 \geqslant C_{2k-1}(2k-2) = C_{2k-2}(2k-2) + I(a_{2k-1} \leqslant 2k-2) = k-1 + I(a_{2k-1} \leqslant 2k-2)$，迫使 $I(a_{2k-1} \leqslant 2k-2) = 0$，即 $a_{2k-1} \geqslant 2k-1$。结合 $a_1=1 \geqslant 1$，对任意 $n \in \mathbb{N}^*$ 均有 $a_n \geqslant n$。

(3) 由(2)推导可知 $C_{2m}(2m-1) = m$ 且 $x_{2m, m} \leqslant 2m-1$，结合 $x_{2m, m} + x_{2m, m+1} = 4m$ 推得 $x_{2m, m+1} \geqslant 2m+1$。这意味着数列中不可能出现偶数，否则若 $a_n = 2m$，由 $a_n \geqslant n$ 必有 $n \leqslant 2m$，导致 $2m \in S_{2m}$，产生矛盾。
因此，当 $k$ 为偶数时，可能解集仅有 $\emptyset$。
当 $k$ 为奇数时，令 $k = 2m-1$。为满足前 $2m$ 项恰有 $m$ 个数 $\leqslant 2m-1$ 且无偶数，这 $m$ 个数必须是 $1, 3, \dots, 2m-1$ 的无重复排列。同理分析 $n=2m-2$ 时的中位数可知 $x_{2m-2, m-1} = 2m-3$，故 $x_{2m-2, m} = 2m-1$。这说明元素 $2m-1$ 必定已出现在 $S_{2m-2}$ 中，其位置 $n \leqslant 2m-2 = k-1$。
结合无重复元素及 $a_n \geqslant n$ 的约束：当 $k=1$ 时，$n \leqslant 1$，解集仅有 $\{1\}$；当 $k=3$ 时，$n \leqslant 2$ 且 $a_1=1$，解集仅有 $\{2\}$。对于奇数 $k \geqslant 5$，因为 $a_1, a_2$ 已被固定为 $1, 3$，新元素 $n$ 的范围放宽至 $3 \leqslant n \leqslant k-1$。在此范围内，总能通过灵活安排后续冗余奇数来构造合法的数列，使得 $a_n=k$ 成立。
综上，关于 $n$ 的方程 $a_n=k$ 所有可能的解集为：当 $k$ 为偶数时为 $\emptyset$；当 $k=1$ 时为 $\{1\}$；当 $k=3$ 时为 $\{2\}$；当奇数 $k \geqslant 5$ 时为 $\{3\}, \{4\}, \dots, \{k-1\}$。
\end{solution}



\begin{example}{邪帝原创导数题}{}
    已知函数$f(x)=x-(a+1)\ln x-\frac{a}{x},(a>1)$\\
    （1）讨论$f(x)$的单调性；\\
    （2）若$f(x_1)=f(x_2)=f(x_3),(x_1<x_2<x_3)$，证明$f(x_1x_2x_3)>1-a$.
\end{example}




\begin{example}{}{}
    已知函数 $f(x)=\ln x+ax^{2}+bx(a,b\in R)$.\\
(1) 若 $a=\frac{1}{e^{4}}, b=-\frac{3}{e^{2}}$, 求函数 $y=f(x)$ 的单调递减区间;\\
(2) 若存在实数 $b$, 使得函数 $f(x)$ 有三个不同的零点 $x_{1}, x_{2}, x_{3}$.求 $a$ 的取值范围;若 $x_{1}, x_{2}, x_{3}$ 成等差数列, 求证: $x_{2}^{2}>e^{3}$.
\end{example}
\begin{solution}
\textbf{（1）}函数 $f(x) = \ln x + ax^2 + bx$ 的定义域为 $(0, +\infty)$。当 $a = \dfrac{1}{e^4},\ b = -\dfrac{3}{e^2}$ 时，
\begin{align*}
f(x) &= \ln x + \frac{1}{e^4}x^2 - \frac{3}{e^2}x, \\
f'(x) &= \frac{1}{x} + \frac{2}{e^4}x - \frac{3}{e^2} = \frac{2x^2 - 3e^2x + e^4}{e^4 x}.
\end{align*}
在定义域内分母恒正，令 $f'(x) \le 0$，即
\[2x^2 - 3e^2x + e^4 \le 0 \implies (2x - e^2)(x - e^2) \le 0.\]
解得 $\dfrac{e^2}{2} \le x \le e^2$。故 $y = f(x)$ 的单调递减区间为 $\left[\dfrac{e^2}{2},\ e^2\right]$。

\textbf{（2）} ① 命题"存在实数 $b$ 使得 $f(x)$ 有三个不同的零点"等价于方程 $\ln x + ax^2 + bx = 0$ 在 $(0, +\infty)$ 上有三个不同的实数根。两边同除以 $x$ 分离参数：
\[-b = \frac{\ln x}{x} + ax.\]
令 $g(x) = \dfrac{\ln x}{x} + ax\ (x > 0)$，问题转化为直线 $y = -b$ 与 $y = g(x)$ 有 3 个不同交点，这要求 $g(x)$ 既有局部极大值又有局部极小值，即 $g'(x) = 0$ 至少有两个不同正根。
\begin{align*}
g'(x) &= \frac{1 - \ln x}{x^2} + a = \frac{ax^2 - \ln x + 1}{x^2}.
\end{align*}
令 $h(x) = ax^2 - \ln x + 1\ (x > 0)$，则 $g'(x)$ 的符号由 $h(x)$ 决定。
\begin{align*}
h'(x) = 2ax - \frac{1}{x} = \frac{2ax^2 - 1}{x}.
\end{align*}
对参数 $a$ 的符号分类讨论如下：

\begin{enumerate}
\item 若 $a \le 0$：对任意 $x > 0$，$2ax^2 \le 0$，故 $h'(x) < 0$ 恒成立。$h(x)$ 严格单调递减，$h(x) = 0$ 至多 1 个根，从而 $g'(x) = 0$ 至多 1 个变号零点，$g(x)$ 至多一次折返，任意水平直线与 $g(x)$ 最多 2 个交点，无法产生三个零点。故 $a \le 0$ 不满足条件。

\item 若 $a > 0$：$h'(x) = 0$ 的唯一驻点为 $x_0 = \dfrac{1}{\sqrt{2a}}$。$h(x)$ 在 $x_0$ 处取得最小值：
\begin{align*}
h_{\min} &= h\!\left(\frac{1}{\sqrt{2a}}\right) = a \cdot \frac{1}{2a} - \ln\frac{1}{\sqrt{2a}} + 1 = \frac{3}{2} + \frac{1}{2}\ln(2a).
\end{align*}
要求 $h(x) = 0$ 有两个不同正根，需 $h_{\min} < 0$，即
\[\frac{3}{2} + \frac{1}{2}\ln(2a) < 0 \implies 2a < e^{-3} \implies 0 < a < \frac{1}{2e^3}.\]
此时 $h(x) \to +\infty$（$x \to 0^+$ 及 $x \to +\infty$），由零点定理，$h(x) = 0$ 恰有两个根 $x_1' < x_0 < x_2'$。$g'(x)$ 在 $(0, x_1')$ 为正，$(x_1', x_2')$ 为负，$(x_2', +\infty)$ 为正，故 $g(x)$ 在 $x_1'$ 处取极大值，$x_2'$ 处取极小值。又 $g(x) \to -\infty$（$x \to 0^+$），$g(x) \to +\infty$（$x \to +\infty$），取 $-b \in \bigl(g(x_2'),\, g(x_1')\bigr)$ 即得 3 个交点。
\end{enumerate}
综上所述，$a$ 的取值范围为 $\left(0,\ \dfrac{1}{2e^3}\right)$。\\
（2）设三个正实数根 $x_1 < x_2 < x_3$ 成等差数列，公差为 $d > 0$，则 $x_1 = x_2 - d$，$x_3 = x_2 + d$，且由 $x_1 > 0$ 得 $0 < d < x_2$。将三个根分别代入 $f(x) = 0$：
\begin{align*}
&\ln(x_2 - d) + a(x_2 - d)^2 + b(x_2 - d) = 0, \tag{1} \\
&\ln x_2 + ax_2^2 + bx_2 = 0, \tag{2} \\
&\ln(x_2 + d) + a(x_2 + d)^2 + b(x_2 + d) = 0. \tag{3}
\end{align*}
取 $(1) + (3) - 2 \times (2)$，消去含 $b$ 的一次项，二次项与对数项分别化简：
\begin{align*}
&\; b(2x_2 - 2x_2) = 0, \\
&a\bigl[(x_2-d)^2 + (x_2+d)^2 - 2x_2^2\bigr] = 2ad^2, \\
&\ln(x_2-d) + \ln(x_2+d) - 2\ln x_2 = \ln\!\left(1 - \frac{d^2}{x_2^2}\right).
\end{align*}
合并得
\[\ln\!\left(1 - \frac{d^2}{x_2^2}\right) + 2ad^2 = 0.\]
令 $t = \dfrac{d^2}{x_2^2} \in (0, 1)$，上式两边同除以 $t$ 可写为
\[2ax_2^2 = \frac{-\ln(1 - t)}{t}.\]
为建立右端下界，构造 $\varphi(t) = -\ln(1 - t) - t$（$t \in [0, 1)$）。求导得
\[\varphi'(t) = \frac{1}{1-t} - 1 = \frac{t}{1-t}.\]
当 $t \in (0, 1)$ 时 $\varphi'(t) > 0$，$\varphi(t)$ 严格单调递增，故 $\varphi(t) > \varphi(0) = 0$，即 $-\ln(1 - t) > t$。两边同除以 $t$ 得 $\dfrac{-\ln(1-t)}{t} > 1$，代入上式得
\[2ax_2^2 > 1 \implies x_2^2 > \frac{1}{2a}.\]
由（1）的结论 $0 < a < \dfrac{1}{2e^3}$，有 $\dfrac{1}{2a} > e^3$，故
\[x_2^2 > \frac{1}{2a} > e^3.\]
即 $x_2^2 > e^3$。
\end{solution}



